% GENERATED FILE. Edit canonical lessons, then run npm run generate:books.
% canonical-source-hash: fnv1a64:e74adb0f41c4a9ba
% canonical-lessons: SA-C257-simantah, SA-C257-lokayanam, SA-C257-karayanam, SA-C257-agamanam, SA-C257-prasthanam, SA-R257-first-pass-words-from-chapters-250-253, SA-R257-second-pass-words-from-chapters-254-257

\chapter{Border, Bus, Car, Arrival, Departure}
\label{ch:sa-border-bus-car-arrival-departure}
\hlchaptermodality{257}

\begin{chapteropening}
\textbf{By the end of this chapter:} \emph{I can say a border, a bus, a car, arrival and departure.}

Complete the last lesson of chapter 257: I can say a border, a bus, a car, arrival and departure.
\end{chapteropening}

\section[\texorpdfstring{\sk{सीमान्तः} (sīmāntaḥ)}{सीमान्तः (sīmāntaḥ)}]{\sk{सीमान्तः} (sīmāntaḥ) — a border}

\label{lesson:SA-C257-simantah}



\begin{quote}
Before the new one: say the Sanskrit for a thought, then the Sanskrit for an opinion.
\end{quote}

\subsection*{You'll want to know: \sk{सीमान्तः}}
\textbf{\sk{सीमान्तः}} — \emph{sīmāntaḥ} — \textquotedblleft{}a border\textquotedblright{}.

\begin{grammarlens}[title={What you've built}]
One more word for this chapter.
\end{grammarlens}

\subsection*{Your turn}
\begin{itemize}
\raggedright
  \item \emph{Say it:} \emph{sīmāntaḥ}
  \item \emph{Say it:} \emph{sīmāntaḥ}, once more
  \item \emph{Say it:} say \emph{sīmāntaḥ}
  \item \emph{Recall:} say the Sanskrit for a thought, then the Sanskrit for an opinion, then say \emph{sīmāntaḥ} again
\end{itemize}

\begin{tcolorbox}[breakable,colback=teal!4,colframe=teal!35!black,title={Before you move on}]
What does \sk{सीमान्तः} mean? (\textquotedblleft{}A border\textquotedblright{}.) Say it once more.
\end{tcolorbox}
\section[\texorpdfstring{\sk{लोकयानम्} (lokayānam)}{लोकयानम् (lokayānam)}]{\sk{लोकयानम्} (lokayānam) — a bus}

\label{lesson:SA-C257-lokayanam}



\begin{quote}
Before the new one: say the Sanskrit for an opinion, then the Sanskrit for a border.
\end{quote}

\subsection*{You'll want to know: \sk{लोकयानम्}}
\textbf{\sk{लोकयानम्}} — \emph{lokayānam} — \textquotedblleft{}a bus\textquotedblright{}.

\begin{grammarlens}[title={What you've built}]
One more word for this chapter.
\end{grammarlens}

\subsection*{Your turn}
\begin{itemize}
\raggedright
  \item \emph{Say it:} \emph{lokayānam}
  \item \emph{Say it:} \emph{lokayānam}, once more
  \item \emph{Say it:} say \emph{lokayānam}
  \item \emph{Recall:} say the Sanskrit for an opinion, then the Sanskrit for a border, then say \emph{lokayānam} again
\end{itemize}

\begin{tcolorbox}[breakable,colback=teal!4,colframe=teal!35!black,title={Before you move on}]
What does \sk{लोकयानम्} mean? (\textquotedblleft{}A bus\textquotedblright{}.) Say it once more.
\end{tcolorbox}
\section[\texorpdfstring{\sk{कारयानम्} (kārayānam)}{कारयानम् (kārayānam)}]{\sk{कारयानम्} (kārayānam) — a car}

\label{lesson:SA-C257-karayanam}



\begin{quote}
Before the new one: say the Sanskrit for a border, then the Sanskrit for a bus.
\end{quote}

\subsection*{You'll want to know: \sk{कारयानम्}}
\textbf{\sk{कारयानम्}} — \emph{kārayānam} — \textquotedblleft{}a car\textquotedblright{}.

\begin{grammarlens}[title={What you've built}]
One more word for this chapter.
\end{grammarlens}

\subsection*{Your turn}
\begin{itemize}
\raggedright
  \item \emph{Say it:} \emph{kārayānam}
  \item \emph{Say it:} \emph{kārayānam}, once more
  \item \emph{Say it:} say \emph{kārayānam}
  \item \emph{Recall:} say the Sanskrit for a border, then the Sanskrit for a bus, then say \emph{kārayānam} again
\end{itemize}

\begin{tcolorbox}[breakable,colback=teal!4,colframe=teal!35!black,title={Before you move on}]
What does \sk{कारयानम्} mean? (\textquotedblleft{}A car\textquotedblright{}.) Say it once more.
\end{tcolorbox}
\section[\texorpdfstring{\sk{आगमनम्} (āgamanam)}{आगमनम् (āgamanam)}]{\sk{आगमनम्} (āgamanam) — arrival}

\label{lesson:SA-C257-agamanam}



\begin{quote}
Before the new one: say the Sanskrit for a bus, then the Sanskrit for a car.
\end{quote}

\subsection*{You'll want to know: \sk{आगमनम्}}
\textbf{\sk{आगमनम्}} — \emph{āgamanam} — \textquotedblleft{}arrival\textquotedblright{}.

\begin{grammarlens}[title={What you've built}]
One more word for this chapter.
\end{grammarlens}

\subsection*{Your turn}
\begin{itemize}
\raggedright
  \item \emph{Say it:} \emph{āgamanam}
  \item \emph{Say it:} \emph{āgamanam}, once more
  \item \emph{Say it:} say \emph{āgamanam}
  \item \emph{Recall:} say the Sanskrit for a bus, then the Sanskrit for a car, then say \emph{āgamanam} again
\end{itemize}

\begin{tcolorbox}[breakable,colback=teal!4,colframe=teal!35!black,title={Before you move on}]
What does \sk{आगमनम्} mean? (\textquotedblleft{}Arrival\textquotedblright{}.) Say it once more.
\end{tcolorbox}
\section[\texorpdfstring{\sk{प्रस्थानम्} (prasthānam)}{प्रस्थानम् (prasthānam)}]{\sk{प्रस्थानम्} (prasthānam) — departure}

\label{lesson:SA-C257-prasthanam}



\begin{quote}
Before the new one: say the Sanskrit for a car, then the Sanskrit for arrival.
\end{quote}

\subsection*{You'll want to know: \sk{प्रस्थानम्}}
\textbf{\sk{प्रस्थानम्}} — \emph{prasthānam} — \textquotedblleft{}departure\textquotedblright{}.

\begin{grammarlens}[title={What you've built}]
One more word for this chapter.
\end{grammarlens}

\subsection*{Your turn}
\begin{itemize}
\raggedright
  \item \emph{Say it:} \emph{prasthānam}
  \item \emph{Say it:} \emph{prasthānam}, once more
  \item \emph{Say it:} say \emph{prasthānam}
  \item \emph{Recall:} say the Sanskrit for a car, then the Sanskrit for arrival, then say \emph{prasthānam} again
\end{itemize}

\begin{tcolorbox}[breakable,colback=teal!4,colframe=teal!35!black,title={Before you move on}]
What does \sk{प्रस्थानम्} mean? (\textquotedblleft{}Departure\textquotedblright{}.) Say it once more.
\end{tcolorbox}
\section[(dialogue)]{Words from chapters 250-253 — a first pass}

\label{lesson:SA-R257-first-pass-words-from-chapters-250-253}



\begin{quote}
Say the words for arrival and departure.
\end{quote}

\subsection*{Your turn}
Say these aloud:

\begin{itemize}
\raggedright
  \item a fort, a palace, a hut, a building, lodging
  \item the earth, the universe, the world, nature, mud
  \item a vine, a thorn, a bud, jasmine, holy basil
  \item an insect, a spider, a scorpion, a lizard, a crocodile
\end{itemize}

\begin{tcolorbox}[breakable,colback=teal!4,colframe=teal!35!black,title={Before you move on}]
A fort? (\textbf{\sk{दुर्गम्}}, \emph{durgam}.) A palace? (\textbf{\sk{प्रासादः}}, \emph{prāsādaḥ}.) A hut? (\textbf{\sk{कुटीरम्}}, \emph{kuṭīram}.) A building? (\textbf{\sk{भवनम्}}, \emph{bhavanam}.) Lodging? (\textbf{\sk{आवासः}}, \emph{āvāsaḥ}.) The earth? (\textbf{\sk{पृथिवी}}, \emph{pṛthivī}.) The universe? (\textbf{\sk{जगत्}}, \emph{jagat}.) The world? (\textbf{\sk{लोकः}}, \emph{lokaḥ}.) Nature? (\textbf{\sk{प्रकृतिः}}, \emph{prakṛtiḥ}.) Mud? (\textbf{\sk{पंकः}}, \emph{paṅkaḥ}.) A vine? (\textbf{\sk{वल्ली}}, \emph{vallī}.) A thorn? (\textbf{\sk{कण्टकः}}, \emph{kaṇṭakaḥ}.) A bud? (\textbf{\sk{कलिका}}, \emph{kalikā}.) Jasmine? (\textbf{\sk{मल्लिका}}, \emph{mallikā}.) Holy basil? (\textbf{\sk{तुलसी}}, \emph{tulasī}.) An insect? (\textbf{\sk{कीटः}}, \emph{kīṭaḥ}.) A spider? (\textbf{\sk{लूता}}, \emph{lūtā}.) A scorpion? (\textbf{\sk{वृश्चिकः}}, \emph{vṛścikaḥ}.) A lizard? (\textbf{\sk{गोधा}}, \emph{godhā}.) A crocodile? (\textbf{\sk{नक्रः}}, \emph{nakraḥ}.)
\end{tcolorbox}
\section[(dialogue)]{Words from chapters 254-257 — a second pass}

\label{lesson:SA-R257-second-pass-words-from-chapters-254-257}



\begin{quote}
Say the words for a ruddy goose and a hawk.
\end{quote}

\subsection*{Your turn}
Say these aloud:

\begin{itemize}
\raggedright
  \item a ruddy goose, a hawk, breakfast, dinner, cooked food
  \item education, knowledge, practice, homework, an answer
  \item training, a conference, a decision, a thought, an opinion
  \item a border, a bus, a car, arrival, departure
\end{itemize}

\begin{tcolorbox}[breakable,colback=teal!4,colframe=teal!35!black,title={Before you move on}]
A ruddy goose? (\textbf{\sk{चक्रवाकः}}, \emph{cakravākaḥ}.) A hawk? (\textbf{\sk{श्येनः}}, \emph{śyenaḥ}.) Breakfast? (\textbf{\sk{प्रातराशः}}, \emph{prātarāśaḥ}.) Dinner? (\textbf{\sk{रात्रिभोजनम्}}, \emph{rātribhojanam}.) Cooked food? (\textbf{\sk{पक्वान्नम्}}, \emph{pakvānnam}.) Education? (\textbf{\sk{शिक्षा}}, \emph{śikṣā}.) Knowledge? (\textbf{\sk{ज्ञानम्}}, \emph{jñānam}.) Practice? (\textbf{\sk{अभ्यासः}}, \emph{abhyāsaḥ}.) Homework? (\textbf{\sk{गृहकार्यम्}}, \emph{gṛhakāryam}.) An answer? (\textbf{\sk{उत्तरम्}}, \emph{uttaram}.) Training? (\textbf{\sk{प्रशिक्षणम्}}, \emph{praśikṣaṇam}.) A conference? (\textbf{\sk{सम्मेलनम्}}, \emph{sammelanam}.) A decision? (\textbf{\sk{निर्णयः}}, \emph{nirṇayaḥ}.) A thought? (\textbf{\sk{विचारः}}, \emph{vicāraḥ}.) An opinion? (\textbf{\sk{मतम्}}, \emph{matam}.) A border? (\textbf{\sk{सीमान्तः}}, \emph{sīmāntaḥ}.) A bus? (\textbf{\sk{लोकयानम्}}, \emph{lokayānam}.) A car? (\textbf{\sk{कारयानम्}}, \emph{kārayānam}.) Arrival? (\textbf{\sk{आगमनम्}}, \emph{āgamanam}.) Departure? (\textbf{\sk{प्रस्थानम्}}, \emph{prasthānam}.)
\end{tcolorbox}
